\documentclass[../../main.tex]{subfiles}
\graphicspath{{\subfix{../../image/}}}

\begin{document}

\section{Bernoulli 试验与直线上的随机游动}\label{section:ch02-Bernoulli试验与随机游动}

\subsection{Bernoulli 概型}

在许多问题中, 我们对试验感兴趣的是试验中某事件 $A$ 是否发生. 例如在产品抽样检查中注意的是抽到废品, 还是抽到好品; 在掷硬币时注意的是出正面还是出反面; 在股票市场中关心的是涨还是跌. 在这类问题中我们可以把事件域取为 $\mathscr{F}=\{\varnothing, A, \bar{A}, \Omega\}$, 并称出现 $A$ 为"成功", 出现 $\bar{A}$ 为"失败". 这种只有两个可能结果的试验称为 Bernoulli 试验.

有些试验的结果不止两个, 譬如, 在电报传输中, 既要传送字母 $A, B, \cdots, Z$ 等, 又要传送其他符号. 但是假如我们所关心的只是字母在传送中所占的百分比, 而不再区别到底是哪一个字母, 则我们可以把出现字母当作是成功, 出现其他符号一律当作是失败, 这时就可以把问题看作 Bernoulli 试验. 同样地, 显像管的寿命可以是不小于 $0$ 的任一数值, 但是有时根据需要, 我们可以把寿命大于 $50\,000$ 小时的显像管当作合格品, 其余都作为次品. 那么, 这类问题还是可以归结为 Bernoulli 试验. 这种例子可以举出不少.

在 Bernoulli 试验中, 首先是要给出下面概率
\begin{align}\label{align:ch02-Bernoulli概率}
P(A)=p,\quad P(\bar{A})=q
\end{align}
显然 $p\geqslant 0$, $q\geqslant 0$, 且 $p+q=1$.

现在考虑重复进行 $n$ 次独立的 Bernoulli 试验, 这里的"重复", 是指在每次试验中事件 $A$, 从而事件 $\bar{A}$ 出现的概率都保持不变. 这种试验称为 $n$ 重 Bernoulli 试验, 记作 $E^n$.

总之, $n$ 重 Bernoulli 试验有下面四个约定：

(i) 每次试验至多出现两个可能结果之一：$A$ 或 $\bar{A}$;

(ii) $A$ 在每次试验中出现的概率 $p$ 保持不变;

(iii) 各次试验相互独立;

(iv) 共进行 $n$ 次试验.

下面先给出 $n$ 重 Bernoulli 试验的概率空间.

$n$ 重 Bernoulli 试验 $E^n$ 的样本点形如：
\begin{align}\label{align:ch02-Bernoulli样本点}
(\hat{A}_1,\hat{A}_2,\cdots,\hat{A}_n)
\end{align}
其中 $\hat{A}_i$ 是 $A_i$ 或 $\bar{A}_i$, 分别表示第 $i$ 次试验中出现 $A$ 或 $\bar{A}$, 显然这种样本点共有 $2^n$ 个, 这是一个有限样本空间.

像对所有有限样本空间一样, 可以把样本点的任意子集作为事件, 构成事件域, 不过在这个场合一般不用明显写出.

为书写方便起见, 在本节中将把样本点\eqref{align:ch02-Bernoulli样本点}简记为 $\hat{A}_1\hat{A}_2\cdots\hat{A}_n$, 例如, $(A_1,A_2,\cdots,A_{n-1},\bar{A}_n)$ 表示前 $n-1$ 次试验均出现事件 $A$ 而第 $n$ 次试验出现事件 $\bar{A}$, 简记为 $A_1A_2\cdots A_{n-1}\bar{A}_n$.

为了给定样本点\eqref{align:ch02-Bernoulli样本点}的概率, 主要看其中 $A$ 或 $\bar{A}$ 出现的次数, 例如若其中有 $l$ 个 $A$, 从而有 $n-l$ 个 $\bar{A}$, 则利用试验的独立性及\eqref{align:ch02-Bernoulli概率}式, 必有
\begin{align}\label{align:ch02-Bernoulli样本点概率}
P(\hat{A}_1\hat{A}_2\cdots\hat{A}_n)&=P(\hat{A}_1)P(\hat{A}_2)\cdots P(\hat{A}_n)\\
&=p^l q^{n-l}\nonumber
\end{align}
特殊地, $P(A_1A_2\cdots A_{n-1}\bar{A}_n)=p^{n-1}q$.

一般事件的概率由它所含样本点的概率求和得到. 这样一来, 我们已对 $n$ 重 Bernoulli 试验给定了概率空间.

有时也需要考虑可列重 Bernoulli 试验 $E^\infty$ 的场合, 这时样本点形如
\begin{align}\label{align:ch02-可列Bernoulli样本点}
(\hat{A}_1,\hat{A}_2,\cdots,\hat{A}_n,\cdots)
\end{align}
其中 $\hat{A}_i$ 仍取 $A_i$ 或 $\bar{A}_i$, 这时样本空间不再有限, 甚至也不可列, 事实上它可与 $[0,1]$ 区间进行一一对应, 对这种试验就不能把样本空间的任意子集都看作事件.

Bernoulli 试验是一种非常重要的概率模型, 它是"在同样条件下进行重复试验"的一种数学模型, 特别在讨论某事件出现的频率时常用这种模型. 历史上, Bernoulli 概型是概率论中最早研究的模型之一, 也是得到最多研究的模型之一, 在理论上具有重要意义, 在我们这门课程中, 一些较为深入的结果也是结合 Bernoulli 概型进行讨论的. 另一方面, 它有着广泛的实际应用, 例如在工业产品质量检查中, 在当代遗传学中它都占有重要地位.

我们将以 Bernoulli 试验为模型探讨机票超售问题. 每个订座旅客当作一次试验, 则他或到时不登机, 记为 $A$, 或到时登机, 记为 $\bar{A}$; $P(A)$ 按过去统计资料取为 $5\%$. 主要的难点在于能否把旅客是否登机看作是独立的, 显然对于购买团体票的旅客作此假定是不适合的, 此外, 大型的交通堵塞等偶然事件也会使这个假定偏离, 不过在一般场合作此假定还是合适的. 全体订座旅客数 $n$ 作为试验总数, 这便构成 Bernoulli 概型.

\subsection{Bernoulli 概型中的一些分布}

下面我们计算 Bernoulli 概型中所出现的一些事件的概率, 这些概率非常重要.

\subsubsection{Bernoulli 分布}

若只进行一次 Bernoulli 试验, 则或是事件 $A$ 出现, 或是事件 $\bar{A}$ 出现, 其概率由\eqref{align:ch02-Bernoulli概率}给出, 称为 Bernoulli 分布, 这是最简单的情况.

\subsubsection{二项分布}

我们来确定 $n$ 重 Bernoulli 试验中事件 $A$ 出现 $k$ 次的概率, 这概率我们记之为 $b(k;n,p)$.

若以 $B_k$ 记 $n$ 重 Bernoulli 试验中事件 $A$ 正好出现 $k$ 次这一事件, 而以 $A_i$ 表示第 $i$ 次试验中出现事件 $A$, 以 $\bar{A}_i$ 表示第 $i$ 次试验中出现 $\bar{A}$, 则
\begin{align}\label{align:ch02-二项分布事件}
B_k=A_1A_2\cdots A_k\bar{A}_{k+1}\cdots\bar{A}_n+\cdots+\bar{A}_1\bar{A}_2\cdots\bar{A}_{n-k}A_{n-k+1}\cdots A_n
\end{align}
右边的每一项表示在某 $k$ 次试验中出现事件 $A$, 在另外 $n-k$ 次试验中出现 $\bar{A}$, 这种项共有 $\binom{n}{k}$ 个, 而且两两互不相容. 由\eqref{align:ch02-Bernoulli样本点概率}可知\eqref{align:ch02-二项分布事件}中右边各项所对应事件的概率均为 $p^k q^{n-k}$, 利用概率的可加性得
\begin{align*}
P(B_k)=\binom{n}{k}p^k q^{n-k}
\end{align*}
即
\begin{align}\label{align:ch02-二项分布}
b(k;n,p)=\binom{n}{k}p^k q^{n-k},\quad k=0,1,2,\cdots,n
\end{align}

注意到 $b(k;n,p)$, $k=0,1,2,\cdots,n$, 是二项式 $(q+ps)^n$ 展开式中 $s^k$ 项的系数, 因此\eqref{align:ch02-二项分布}称为二项分布（binomial distribution）. 特别地
\begin{align}\label{align:ch02-二项分布归一化}
\sum_{k=0}^{n} b(k;n,p)=\sum_{k=0}^{n} \binom{n}{k}p^k q^{n-k}=(q+p)^n=1
\end{align}

\begin{example}[有放回抽样中的废品数]\label{example:ch02-有放回抽样废品数}
若在 $N$ 件产品中有 $M$ 件废品, 现进行 $n$ 次有放回的抽样检查, 问共抽得 $k$ 件废品的概率是多少？

\begin{proof}[解]
由于抽样是有放回的, 因此这是 $n$ 重 Bernoulli 试验, 若以 $A$ 记各次试验中出现废品这一事件, 则
\begin{align*}
p=P(A)=\frac{M}{N}
\end{align*}
因此所求的概率为
\begin{align*}
b\left(k;n,\frac{M}{N}\right)=\binom{n}{k}\left(\frac{M}{N}\right)^k\left(1-\frac{M}{N}\right)^{n-k}
\end{align*}
这概率在第一章 \S3 中曾出现过.
\end{proof}
\end{example}

\begin{example}[群体遗传学中的基因频率]\label{example:ch02-群体遗传学基因频率}
在群体遗传学中, 假定可遗传的指标是依赖于基因的. 基因总是成对出现并且具有两种形式 $A$ 及 $a$ 中的一种. 假定每一代具有 $2N$ 个基因, 则其中 $A$ 所占的比例数 $\dfrac{i}{2N}$ 称为基因频率. 基因频率的变化过程与群体的进化情况有着密切的关系.

如果进行的是随机交配, 也就是说任何个体有同样的机会和任何其他个体配种, 则遗传学中对基因的遗传作下面假定：子代个体是按 Bernoulli 概型从上一代每个亲体中取得基因的.

因此若上一代的基因频率为 $\dfrac{i}{2N}$, 则下一代的基因频率为 $\dfrac{j}{2N}$ 的概率由二项分布给出
\begin{align*}
p_{ij}=\binom{2N}{j}\left(\frac{i}{2N}\right)^j\left(1-\frac{i}{2N}\right)^{2N-j}
\end{align*}

在这个例子中基因频率 $\dfrac{i}{2N}$ 相当于 Bernoulli 试验中出现成功的概率 $p$.
\end{example}

\subsubsection{几何分布}

现在讨论在 Bernoulli 试验中首次成功出现在第 $k$ 次试验的概率, 要使首次成功出现在第 $k$ 次试验, 必须而且只需在前 $k-1$ 次试验中都出现事件 $\bar{A}$, 而第 $k$ 次试验出现 $A$, 因此这事件（记为 $W_k$）可表示为
\begin{align}\label{align:ch02-几何分布事件}
W_k=\bar{A}_1\bar{A}_2\cdots\bar{A}_{k-1}A_k
\end{align}
利用试验的独立性, 其概率为
\begin{align*}
P(W_k)=P(\bar{A}_1)P(\bar{A}_2)\cdots P(\bar{A}_{k-1})P(A_k)=q^{k-1}p
\end{align*}
记
\begin{align}\label{align:ch02-几何分布}
g(k;p)=q^{k-1}p,\quad k=1,2,\cdots
\end{align}
$g(k;p)$ 是几何级数的一般项, 因此\eqref{align:ch02-几何分布}称为几何分布. 这里有
\begin{align}\label{align:ch02-几何分布归一化}
\sum_{k=1}^{\infty} g(k;p)=\sum_{k=1}^{\infty} q^{k-1}p=p\frac{1}{1-q}=1
\end{align}

几何分布给出了等待事件 $A$ 出现共试验了 $k$ 次的概率, 这类概率在许多问题中出现, 我们在第一章 \S3 中曾遇到它. 下面是一个模型化了的例子.

\begin{example}[试钥匙开门]\label{example:ch02-试钥匙开门}
一个人要开门, 他共有 $n$ 把钥匙, 其中仅有一把是能开这门的. 他随机地选取一把钥匙开门, 即在每次试开时每一把钥匙都以概率 $\dfrac{1}{n}$ 被使用, 这人在第 $s$ 次试开时才首次成功的概率是多少？

\begin{proof}[解]
这是一个 Bernoulli 试验, $p=\dfrac{1}{n}$, 由\eqref{align:ch02-几何分布}, 所求概率为
\begin{align*}
g\left(s;\frac{1}{n}\right)=\left(\frac{n-1}{n}\right)^{s-1}\frac{1}{n}
\end{align*}
\end{proof}

应当指出, 讨论事件 $A$ 的首次出现事实上牵涉到可列次 Bernoulli 试验, 因此它的样本空间应取为\eqref{align:ch02-可列Bernoulli样本点}所表示的样本点全体. 这个空间不是可列的, 所以不能把它的一切子集都作为事件. 在上述讨论中, 我们只是把用\eqref{align:ch02-几何分布事件}表出的 $W_k$ 作为事件, 并以\eqref{align:ch02-几何分布}给出其概率, 这里的 $W_k$ 事实上也包含了不可列个用\eqref{align:ch02-可列Bernoulli样本点}表出的样本点. 当然在讨论这个问题时, 我们也可以干脆把 $W_k$ 作为样本点, 这时样本空间取为 $\{W_1, W_2, W_3, \cdots\}$, 并以\eqref{align:ch02-几何分布}给出其概率, 这就成了一个离散样本空间.
\end{example}

下面讨论的是更复杂一点的情况, 即 Pascal 分布, 它可以看作几何分布的一种推广.

\subsubsection{Pascal 分布}

考虑 Bernoulli 试验, 让我们考察要多长时间才会出现第 $r$ 次成功.

若第 $r$ 次成功发生在第 $\zeta$ 次试验, 则必然有 $\zeta\geqslant r$.

让我们以 $C_k$ 表示第 $r$ 次成功发生在第 $k$ 次试验这一事件, 并以 $f(k;r,p)$ 记其概率, $C_k$ 发生当且仅当前面的 $k-1$ 次试验中有 $r-1$ 次成功, $k-r$ 次失败, 而第 $k$ 次试验的结果为成功, 这两个事件的概率分别为 $\binom{k-1}{r-1}p^{r-1}q^{k-r}$ 与 $p$, 于是利用试验的独立性, 得到
\begin{align*}
P(C_k)=\binom{k-1}{r-1}p^{r-1}q^{k-r}\cdot p=\binom{k-1}{r-1}p^r q^{k-r}
\end{align*}
即
\begin{align}\label{align:ch02-Pascal分布}
f(k;r,p)=\binom{k-1}{r-1}p^r q^{k-r},\quad k=r,r+1,\cdots
\end{align}

注意到
\begin{align*}
\sum_{k=r}^{\infty} f(k;r,p)&=\sum_{k=r}^{\infty}\binom{k-1}{r-1}p^r q^{k-r}\\
&=\sum_{l=0}^{\infty}\binom{r+l-1}{r-1}p^r q^l=\sum_{l=0}^{\infty}\binom{r+l-1}{l}p^r q^l\\
&=\sum_{l=0}^{\infty}\binom{-r}{l}(-1)^l p^r q^l=p^r(1-q)^{-r}=1
\end{align*}
这里利用了推广的二项系数公式和牛顿二项式.

$f(k;r,p)$ 称为 Pascal 分布. 特别当 $r=1$ 时, 我们得到几何分布.

Pascal 分布与著名的分赌注问题有关. 这个问题来源甚古, 但直到 de Méré 向 Pascal 提出并导致 Pascal 与 Fermat 的通信, 后又引起 Huygens（1629—1695）的兴趣, 才由他们三人分别给出正确答案. Pascal 和 Fermat 在解题中归结到取胜的概率, 而 Huygens 则引入数学期望的概念. 可以说这宣告了概率论这一学科的诞生.

分赌注问题大意如下：甲、乙两个赌徒按某种方式下注赌博, 说定先胜 $t$ 局者将赢得全部赌注, 但进行到甲胜 $r$ 局, 乙胜 $s$ 局（$r<t$, $s<t$）时, 因故不得不中止, 试问如何分配这些赌注才公平合理？

有人建议用已胜局数作比例分配赌注, 即以 $r:s$ 来分配, 但这种分法显然没有考虑到最终取胜的概率. 若以 $n=t-r$ 及 $m=t-s$ 分别记甲及乙为达到最后胜利所须再胜的局数, 又设甲在每局中取胜的概率为 $p$, 我们便可以把分赌注问题归结为如下概率问题：在 Bernoulli 试验中, 求在出现 $m$ 次 $\bar{A}$ 之前出现 $n$ 次 $A$ 的概率.

若以 $p_{\text{甲}}$ 记上述概率, 则它为甲最终取胜的概率, 那么赌注以 $p_{\text{甲}}:1-p_{\text{甲}}$ 分配是公平合理的. Pascal 和 Fermat 都在某种程度上达到这个结果.

现在, 若利用 Pascal 分布, 则容易写出答案
\begin{align}\label{align:ch02-分赌注答案1}
p_{\text{甲}}=\sum_{k=0}^{m-1}\binom{n+k-1}{k}p^n q^k
\end{align}
或
\begin{align}\label{align:ch02-分赌注答案2}
p_{\text{甲}}=\sum_{k=n}^{\infty}\binom{m+k-1}{k}p^k q^m
\end{align}

另外, 容易证明, 再赌 $n+m-1$ 局一定可以决定胜负. 因此甲为取得最终胜利只须而且必须在后继的 $n+m-1$ 局中至少胜 $n$ 局. 这样, 利用二项分布可以知道,
\begin{align}\label{align:ch02-分赌注答案3}
p_{\text{甲}}=\sum_{k=n}^{n+m-1}\binom{n+m-1}{k}p^k q^{n+m-1-k}
\end{align}
可以证明上述三个答案是一致的（习题 29）.

下面是与 Pascal 分布有关的另一个有名例子.

\begin{example}[Banach 火柴盒问题]\label{example:ch02-Banach火柴盒}
数学家的左、右衣袋中各放有一盒装有 $N$ 根火柴的火柴盒, 每次抽烟时任取一盒用一根, 求发现一盒用光时, 另一盒有 $r$ 根的概率.

\begin{proof}[解]
看作 $p=\dfrac{1}{2}$ 的 Bernoulli 试验. 要左边空而右边剩 $r$ 根, 应该是左边摸过 $N+1$ 次（前 $N$ 次用去 $N$ 根火柴, 最后一次发觉火柴盒是空的）, 而右边摸过 $N-r$ 次, 这事件的概率为
\begin{align*}
f\left(2N-r+1;N+1,\frac{1}{2}\right)=\binom{2N-r}{N}\left(\frac{1}{2}\right)^{2N-r+1}
\end{align*}
对于右边先空的情况可同样考虑, 因此所求的概率为
\begin{align*}
u_r=2\cdot f\left(2N-r+1;N+1,\frac{1}{2}\right)=\binom{2N-r}{N}2^{-2N+r}
\end{align*}
\end{proof}
\end{example}

\subsection{直线上的随机游动}

考虑 $x$ 轴上的一个质点, 假定它只能位于整数点, 在时刻 $t=0$ 时, 它处于初始位置 $a$（$a$ 是整数）, 以后每隔单位时间, 它总受到一个外力的随机作用, 使位置发生变化, 分别以概率 $p$ 及概率 $q=1-p$ 向正的或负的方向移动一个单位, 我们所关心的是质点在时刻 $t=n$ 时的位置. 用这种方式描述的质点运动称为随机游动.

若质点可以在整个数轴的整数点上游动, 则称这种随机游动为无限制随机游动. 若在某点 $d$ 设有一个吸收壁, 质点一到达这点即被吸收而不再游动, 因而整个游动也就结束了, 这种随机游动称为在 $d$ 点有吸收壁的随机游动. 此外还可以考虑带有反射壁及弹性壁的随机游动. 在一个随机游动中还可以具有不止一个壁.

当 $p=q=\dfrac{1}{2}$ 时, 随机游动称为对称的, 这时质点向左或向右移动的可能性相等.

自然科学中的大量问题归结为随机游动问题, 例如随机游动模型可以作为布朗运动的初步近似. 股票价格涨落和汇率变化是否具有随机游动特征, 更是现代金融界最充满火药味的论争之一. 概率论中的一些古典问题也引导到随机游动问题, 事实上, 随机游动可以看作是 Bernoulli 试验的一种描述法.

关于随机游动, 已进行过许多研究, 我们只介绍它的两个最简单的模型.

\subsubsection{无限制随机游动}

假定质点在时刻 $0$ 从原点出发, 以 $S_n$ 记它在时刻 $t=n$ 时的位置. 为了使质点在时刻 $t=n$ 时位于 $k$（$k$ 也可以是负整数）, 必须而且只须在前 $n$ 次游动中向右游动的次数比向左游动的次数多 $k$ 次, 若以 $x$ 记它在前 $n$ 次游动中向右游动的次数, $y$ 记向左移动的次数, 则
\begin{align*}
\begin{cases}
x+y=n\\
x-y=k
\end{cases}
\end{align*}
即 $x=\dfrac{n+k}{2}$, 因为 $x$ 是整数, 所以 $k$ 必须与 $n$ 具有相同的奇偶性.

事件 $\{S_n=k\}$ 发生相当于要求在前 $n$ 次游动中有 $\dfrac{n+k}{2}$ 次向右, $\dfrac{n-k}{2}$ 次向左, 利用二项分布即得
\begin{align*}
P\{S_n=k\}=\binom{n}{\frac{n+k}{2}}q^{\frac{n-k}{2}}p^{\frac{n+k}{2}}
\end{align*}
当 $k$ 与 $n$ 奇偶性相反时, 概率为 $0$.

\subsubsection{两端带有吸收壁的随机游动}

假定质点在时刻 $t=0$ 时, 位于 $x=a$, 而在 $x=0$ 及 $x=a+b$ 处各有一个吸收壁, 我们来求质点在 $x=0$ 被吸收或在 $x=a+b$ 被吸收的概率. 用的是差分方程法.

若以 $q_n$ 记质点的初始位置为 $n$ 而最终在 $a+b$ 点被吸收的概率. 显然
\begin{align}\label{align:ch02-吸收壁边界条件}
q_0=0,\quad q_{a+b}=1
\end{align}

如果某时刻质点位于 $x=n$, 这里 $1\leqslant n\leqslant a+b-1$, 则它要被 $x=a+b$ 吸收, 有两种方式来实现：一种是接下去一次移动是向右的而最终被 $x=a+b$ 吸收; 另一种是接下去一次移动是向左的而最终被 $x=a+b$ 吸收. 所以按全概率公式有
\begin{align}\label{align:ch02-吸收壁差分方程}
q_n=pq_{n+1}+qq_{n-1},\quad n=1,2,\cdots,a+b-1
\end{align}

这样, 我们得到了关于 $q_n$ 的一个二阶差分方程\eqref{align:ch02-吸收壁差分方程}, 再用边界条件\eqref{align:ch02-吸收壁边界条件}就可以求解. 利用这个差分方程系数的特殊性, 比较方便的解法是把\eqref{align:ch02-吸收壁差分方程}改写成
\begin{align*}
p(q_{n+1}-q_n)=q(q_n-q_{n-1}),\quad n=1,2,\cdots,a+b-1
\end{align*}

若记 $c_n=q_{n+1}-q_n$, $r=\dfrac{q}{p}$, 则又能写成
\begin{align*}
c_n=rc_{n-1},\quad n=1,2,\cdots,a+b-1
\end{align*}

下面分两种情况求解：

(i) $r=1$, 即 $p=q=\dfrac{1}{2}$, 也即对称随机游动的场合. 这时 $c_n=c_{n-1}$, 因此, 若记
\begin{align*}
q_{n+1}-q_n=q_n-q_{n-1}=\cdots=q_1-q_0=d
\end{align*}
则
\begin{align*}
q_n=q_0+nd
\end{align*}
由于 $q_0=0$, $q_{a+b}=1$, 故有
\begin{align*}
q_n=\frac{n}{a+b}
\end{align*}
特别地
\begin{align}\label{align:ch02-对称吸收壁结果}
q_a=\frac{a}{a+b}
\end{align}

(ii) $r\neq 1$, 即 $p\neq q$ 的场合.

这时
\begin{align*}
c_n=rc_{n-1}=r^2c_{n-2}=\cdots=r^n c_0
\end{align*}
从而
\begin{align*}
q_n-q_0=\sum_{k=0}^{n-1}(q_{k+1}-q_k)=\sum_{k=0}^{n-1} c_k=\sum_{k=0}^{n-1} r^k c_0=\frac{1-r^n}{1-r}c_0
\end{align*}
由于 $q_0=0$, $q_{a+b}=1$, 故有
\begin{align*}
\frac{1-r^{a+b}}{1-r}c_0=1
\end{align*}
因此
\begin{align*}
q_n=\frac{1-r^n}{1-r^{a+b}}
\end{align*}
特别地
\begin{align}\label{align:ch02-非对称吸收壁结果}
q_a=\frac{1-r^a}{1-r^{a+b}}=\frac{1-\left(\frac{q}{p}\right)^a}{1-\left(\frac{q}{p}\right)^{a+b}}
\end{align}

若以 $p_n$ 记质点自 $n$ 出发而在 $0$ 点被吸收的概率, 同样可以列出差分方程
\begin{align*}
p_n=pp_{n+1}+qp_{n-1},\quad n=1,2,\cdots,a+b-1
\end{align*}
及边界条件
\begin{align*}
p_0=1,\quad p_{a+b}=0
\end{align*}

类似地可以求得, 当 $p=q=\dfrac{1}{2}$ 时
\begin{align}\label{align:ch02-对称吸收壁p结果}
p_a=\frac{b}{a+b}
\end{align}
而在 $p\neq q$ 的场合
\begin{align}\label{align:ch02-非对称吸收壁p结果}
p_a=\frac{1-\left(\frac{p}{q}\right)^b}{1-\left(\frac{p}{q}\right)^{a+b}}=\frac{\left(\frac{q}{p}\right)^a-\left(\frac{q}{p}\right)^{a+b}}{1-\left(\frac{q}{p}\right)^{a+b}}
\end{align}

不管在什么场合, 都有
\begin{align*}
p_a+q_a=1
\end{align*}
也就是说随机游动的质点最终一定要被两个端点之一所吸收.

注意, \eqref{align:ch02-对称吸收壁结果} 及 \eqref{align:ch02-对称吸收壁p结果} 也可以通过 \eqref{align:ch02-非对称吸收壁结果} 及 \eqref{align:ch02-非对称吸收壁p结果}, 令 $p\to q$ 用洛必达法则得到.

顺便指出, 两端有吸收壁的随机游动与概率论发展史上有名的赌徒输光问题有密切联系. 这个问题是这样叙述的：甲、乙进行赌博, 其赌本分别为 $a$ 及 $b$, 若每局赌注为 $1$, 而甲、乙在每局中赢的概率分别为 $p$ 及 $q$, 试求乙（或甲）把赌本输光的概率.

这个问题最初由 Huygens 对公平博弯即 $p=q=\dfrac{1}{2}$ 的情况作了讨论, 后由 Bernoulli 推广到一般场合.

\subsection{推广的 Bernoulli 试验与多项分布}

二项分布可以容易地推广到 $n$ 次重复独立试验且每次试验可能有若干个结果的情形. 把每次试验的可能结果记为 $A_1, A_2, \cdots, A_r$, 而 $P(A_i)=p_i$, $i=1,2,\cdots,r$, 且
\begin{align}\label{align:ch02-多项分布概率和}
p_1+p_2+\cdots+p_r=1,\quad p_i\geqslant 0
\end{align}
当 $r=2$ 时, 我们得到 Bernoulli 试验.

在这种推广的 Bernoulli 试验中, 不难导出, 在 $n$ 次试验中 $A_1$ 出现 $k_1$ 次, $A_2$ 出现 $k_2$ 次, $\cdots\cdots$ $A_r$ 出现 $k_r$ 次的概率为
\begin{align}\label{align:ch02-多项分布}
\frac{n!}{k_1!\ k_2!\ \cdots\ k_r!}p_1^{k_1}p_2^{k_2}\cdots p_r^{k_r}
\end{align}
这里 $k_i\geqslant 0$, 且 $k_1+k_2+\cdots+k_r=n$.

公式\eqref{align:ch02-多项分布}称为多项分布, 因为它是 $(p_1+p_2+\cdots+p_r)^n$ 的展开式的一般项, 而且由\eqref{align:ch02-多项分布概率和}知
\begin{align*}
\sum_{\substack{k_i\geqslant 0\\ k_1+k_2+\cdots+k_r=n}}\frac{n!}{k_1!\ k_2!\ \cdots\ k_r!}p_1^{k_1}p_2^{k_2}\cdots p_r^{k_r}=1
\end{align*}

显然多项分布是二项分布的推广. 二项分布中的许多结果都能平行地推广到多项分布的场合, 以后我们只详细讨论二项分布的有关问题.

在产品检查中, 若对产品质量所用的标准不只是好品与废品, 而是分得更细, 例如有一等品, 二等品, 三等品, 等外品四类, 则从中取出 $n$ 件, 求一等品有 $k_1$ 件, 二等品有 $k_2$ 件, 三等品有 $k_3$ 件, 等外品有 $k_4$ 件的概率时便得到多项分布.

\begin{example}[人类血型分布]\label{example:ch02-人类血型分布}
人类的血型分为 $O$, $A$, $B$, $AB$ 四型, 假定某地区的居民中这四种血型的人的百分比分别为 $0.4$, $0.3$, $0.25$, $0.05$, 若从此地区居民中随机地选出 $5$ 人, 求有两个为 $O$ 型, 其他三个分别是 $A$, $B$, $AB$ 型的概率.

\begin{proof}[解]
推广的 Bernoulli 试验可以用于这个场合, 所求的概率为
\begin{align*}
P=\frac{5!}{2!\ 1!\ 1!\ 1!}\times 0.4^2\times 0.3\times 0.25\times 0.05=0.036
\end{align*}
\end{proof}
\end{example}

也可以研究平面上或空间的随机游动, 下面是简单的一例.

\begin{example}[平面上随机游动的返回概率]\label{example:ch02-平面随机游动返回}
一质点从平面上某点出发, 等可能地向上, 下, 左, 右方向移动, 每次移动的距离为 $1$, 求经过 $2n$ 次移动后回到出发点的概率.

\begin{proof}[解]
这可以归结为上述推广的 Bernoulli 试验的问题, 分别以事件 $A_1, A_2, A_3, A_4$ 表示质点向上, 下, 左, 右移动一格, 则 $p_1=p_2=p_3=p_4=\dfrac{1}{4}$. 若要在 $2n$ 次移动后回到原来的出发点, 则向左移动的次数与向右移动的次数应该相等, 向上移动的次数与向下移动的次数也应该相等. 而总移动次数为 $2n$, 因此所求概率为
\begin{align*}
P&=\sum_{k+m=n}\frac{(2n)!}{(k!)^2(m!)^2}\left(\frac{1}{4}\right)^{2n}\\
&=\sum_{k=0}^{n}\frac{(2n)!}{(k!)^2[(n-k)!]^2}\left(\frac{1}{4}\right)^{2n}\\
&=\left(\frac{1}{4}\right)^{2n}\frac{(2n)!}{(n!)^2}\sum_{k=0}^{n}\left[\frac{n!}{k!\ (n-k)!}\right]^2\\
&=\left(\frac{1}{4}\right)^{2n}\binom{2n}{n}\sum_{k=0}^{n}\binom{n}{k}^2=\left(\frac{1}{4}\right)^{2n}\binom{2n}{n}^2
\end{align*}
最后一个等式用到第一章的恒等式.
\end{proof}
\end{example}

\end{document}